\section{Targeted Skill Intervention}
\label{sec:intervention}
Each frozen endpoint is replayed under three conditions while its evidence package, model configuration, prompt harness, evaluator, and period remain fixed. The intervention is intentionally simple: compare the original agent with the same agent receiving either a target-blind helper or the one reusable operation predicted by the registered failure mechanism.

\paragraph{Design requirements.}
The comparison must satisfy three requirements. First, the targeted arm must improve over the original no-skill agent; otherwise there is no repaired failure. Second, it must also improve over a helper with the same callable surface and matched implementation complexity; otherwise the effect may reflect generic assistance rather than the target operation. Third, the endpoint outcome and transfer split must be frozen before the corresponding model outcomes are opened, so neither the mechanism label nor the success criterion can move to fit the result.

The targeted condition injects the registered mechanism-specific procedure. The generic condition injects a frozen target-blind helper under the same callable wrapper. The no-skill condition adds no experimental helper. The scientific unit is the frozen endpoint, not a repeat. We evaluate the two contrasts from Definition~1 jointly: targeted minus no skill asks whether the operation repairs the original agent, while targeted minus generic asks whether that repair exceeds generic helper/code assistance. Neither contrast is sufficient alone.

\subsection{Mechanism contracts}
The cutoff skill filters candidate evidence by the target cutoff. Its generic control validates package structure without using timestamps. The release-alignment skill normalizes the reporting period represented by a source release. Its generic control inspects document structure without performing date-to-period mapping. The grounding skill enriches candidate documentary spans using reusable attribution and outlook cues. Its control can expose document structure but has no contextual event selector.

No reusable skill contains scenario-specific values, dates, labels, or expected answers. Skill source hashes are frozen before execution. The targeted and generic source implementations are matched exactly in lexical-token count and AST-node count for each family. This generic control is deliberately \emph{target-blind}, not assumed behaviorally neutral: it controls the callable wrapper and implementation-complexity surface while withholding the temporal operation under test. The no-skill arm is therefore essential for detecting harm introduced by the particular generic helper.

\subsection{Resource parity and use semantics}
The experiment measures the effect of controlled procedure exposure, not autonomous skill selection. Under the frozen harness, a targeted or generic helper is executed once before the model answer whenever that arm is assigned; the no-skill arm executes no experimental helper. Targeted and generic arms have equal external tool-call budgets and matched source complexity. Consequently, our efficiency lane is a parity audit rather than a speedup claim: we test whether the mechanism-specific operation changes outcomes under comparable helper budget. We do not claim lower wall-clock latency or lower provider cost. Post-hoc token-length and condition-position diagnostics are reported separately in Section~\ref{sec:results} and Appendix~A.

\subsection{Frozen family outcomes and scorer timeline}
Before any source-native model outcome is opened, the failure-family contracts define what constitutes success. A cutoff success requires a nonempty set of cited evidence and a selected latest reference that are all available by the target cutoff. A release-alignment success requires the correct reporting-period identity together with the aligned headline value. A grounding success requires the requested evidence slots to be filled by frozen first-party sentence IDs in the correct slot order.

An auxiliary full-task conjunctive scorer is preserved only as sensitivity analysis because it adds orthogonal output fields and exact surface-form requirements beyond the registered mechanism outcome. The family definitions frozen before execution remain primary; Appendix~A records the scorer timeline and representation audit.

The primary decision rule follows Definition~1. A positive targeted-minus-generic contrast with a nonpositive targeted-minus-no-skill contrast is not counted as repair. This pattern indicates that the generic helper may have degraded behavior while the targeted condition merely restored the original baseline. Conversely, targeted-minus-no-skill alone cannot distinguish mechanism-specific repair from a generic benefit of receiving an additional helper. The bottleneck label therefore requires both comparisons on the same endpoint family.

\paragraph{Treatment-surface boundary.}
The experiment intervenes on an operation packaged as callable agent state. It does not compare that package against a behaviorally equivalent timestamp-aware retriever or retrieval-side temporal prefilter. Such a treatment would answer a different estimand---whether the \emph{container} matters once the temporal operation is held behaviorally equivalent. We leave that comparison unexecuted and make no container-superiority claim.

\subsection{Held-out transfer and evidence layers}
The initial replay freezes twelve pre-April endpoints as C3-R and eleven April--May endpoints as held-out C4-R; same-domain and compatible cross-domain labels are fixed before outcomes, and skill source is unchanged across the split. A separate EIA validation freezes four earlier and eight held-out energy endpoints while reusing T1/G1 byte-identically. DeepSeek is the registered primary track, Kimi a provenance-clean replication selected by a pre-outcome support rule, and the mini-model layer an exploratory capacity stress. Full execution chronology, the failed local-secondary gate, and row-level provenance bindings are reported in Appendix~A.
